As faixas de frequência auditíva aproximadas, com base em dados da literatura, são:

\begin{itemize}
\item \textbf{Humanos:} 20 Hz a 20\,000 Hz (20 kHz). Faixa padrão para adultos jovens.
\item \textbf{Elefantes (africanos e asiáticos):} $\approx$14 Hz a 12\,000 Hz. A faixa se estende ao infrassom (abaixo de 20 Hz), permitindo comunicação de longa distância por sons com frequências de 14--35 Hz.
\item \textbf{Baleias azuis:} $\approx$10 Hz a 31\,000 Hz. Produzem sons de baixa frequência (10--40 Hz) propagados por longas distâncias no oceano.
\item \textbf{Baleias jubarte (corcunda):} $\approx$20 Hz a 24\,000 Hz.
\end{itemize}

\textbf{Comparação:}
\begin{itemize}
\item Elefantes e baleias cobrem frequências \textbf{abaixo de 20 Hz} (infrassom), inacessíveis à audição humana.
\item A relação de escala é física: animais maiores têm vias vocáis maiores, produzindo ondas de comprimento maior e, portanto, frequência menor.
\item A frequência mínima auditíva escala aproximadamente como $f_{\min} \propto L^{-1}$, onde $L$ é o comprimento característico do animal.
\end{itemize}

